\section{典型国产算力推理引擎与适配路径}

如果沿用前文“一个底座、三条主路径”的分析框架来观察本土路线分化，可以更清楚地看出不同方案真正分开的并不只是产品形态，而是它们把复杂度压在何处：有的路线优先掌握执行与通信主路径，以最快速度把国产后端接入主干框架；有的路线优先把状态治理、部署流程和平台软件栈整体收拢，以换取长期交付确定性；还有的路线试图在统一接口不破坏的前提下，同时经营执行优化、状态治理与压缩成本协同三条主路径。也正因此，比较不同国产推理引擎时，真正需要问的问题不只是“兼容谁”或“性能上限如何”，而是其是否具备稳定的统一抽象和观测底座，以及它把执行、状态和成本三个关键矛盾分别如何组织进系统中。

进一步说，路线比较不能只停留在抽象层。真正决定系统成熟度的，是这些路线是否显式处理了统一执行 IR、阶段解耦与通信运行时，是否提供了前缀索引、多级驻留和状态迁移能力，是否把混合精度、KV 压缩和单位 token 成本带入同一优化闭环。只有把这些技术点阐明，才能保证 survey 不会漏掉真正决定国产推理系统上限的技术问题。

\subsection{面向通用开源引擎的国产后端适配}

以 vLLM 与 SGLang 为代表的开源推理框架，正在通过平台接口、后端插件、模型执行器扩展和定制算子等方式适配国产硬件\citep{vllm,sglang,sglangpaper,vllmascend}。该路线的优势在于可以持续复用上游社区在调度、缓存、结构化输出与多模态支持方面的能力演进，并借助社区已有的模型生态快速覆盖主流开源模型。这一路线的核心思想不是重写一个本土专用引擎，而是在保留上游主干能力的同时，为国产后端补齐执行和运维能力。典型实现方式包括：在平台层增加设备选择与能力注册；在 worker 或 model runner 层注入后端特化逻辑；在注意力、归一化、采样等热点路径增加定制算子；以及在调度与缓存层引入面向本土平台的约束门控。其好处是迭代快、模型覆盖广、与上游社区距离较近；不足则在于，若抽象边界设计不当，后端 patch 会随着上游快速演进不断累积技术债。在国内部署实践中，这一路线的真实工作量通常还不止于“补齐 kernel 或后端接口”。一方面，系统往往需要同时处理模型接入、量化转换、OpenAI 兼容协议、监控埋点和部署脚本等外围工程；另一方面，后端适配是否成功，往往还取决于模型支持矩阵更新速度、异常回退路径是否可观测，以及版本升级时能否继续跟上上游抽象变化。LMDeploy 更强调部署工具链与模型接入效率，vLLM-Ascend 更能体现对上游执行边界的依赖，而 MindIE 一类方案则把这些问题更多吸收到平台交付链路内部\citep{lmdeploy,vllmascend,mindie}。因此，本土团队在选择“开源主干后端适配”时，实际评估的往往不是单一性能收益，而是模型跟进速度、改动侵入度和后续运维成本之间的综合平衡。若进一步看国内 fork 的工程组织方式，vLLM-HUST 更适合被视作介于“开源主干后端适配”和“混合式架构”之间、但整体更偏混合式的一类实践：它在上层继续保持 vLLM 兼容的服务接口、多模态、Prefix Cache 与 Multi-LoRA 等能力边界，在下层则通过平台插件、Ascend 运行时预检和外部运行时治理工具，把国产平台接入与环境治理从核心执行链中拆开\citep{vllmhust}。这意味着它并不是单纯在主干里补一个后端，而是在尽量复用上游 serving 主体的前提下，把平台特化与启动护栏做成相对独立、可演进的工程层。这一类实践对综述尤其重要，因为它揭示了国产项目在“复用上游”之外的另一层真实工作：不是简单选择跟或不跟上游，而是重新设计哪些能力必须留在主干、哪些能力应沉淀为插件、哪些治理逻辑需要放在运行时外围工具链中。如果说 vLLM-Ascend 更像把国产平台能力直接压入主干兼容边界，那么 vLLM-HUST 则更强调在主干周围补齐平台护栏、能力注册和环境治理。这种差异说明，开源主干后端适配并不是单一形态，它内部也在分化为“主链后端实现优先”和“主链复用 + 外围治理增强”两种组织方向。

更具体地看，vLLM-Ascend 与 vLLM-MLU 应被视为同一技术族：二者都尽量继承 vLLM 的 OpenAI 兼容接口、调度语义和缓存抽象，再分别把 Ascend NPU 与寒武纪 MLU 的执行、依赖和分布式适配压入 plugin backend 边界\citep{vllmascend,vllmmlu}。SGLang Ascend 则不是另一个独立国产引擎，而是 SGLang 主线仓库中的 Ascend 平台后端；其意义在于 RadixAttention、PD 分离、structured outputs、chunked prefill 等主流 serving 机制正在被映射到国产 NPU 后端，而不是在本土分支里孤立重做\citep{sglang,sglang_ascend_npu}。这一路线的证据边界也应写清：插件式后端的成熟度受 CANN、torch\_npu、厂商 SDK、容器镜像和上游版本节奏共同约束；SGLang Ascend 中 NPUGraph、HiCache/MoonCake、量化和投机解码等更高阶能力仍有不少属于 roadmap 或持续工程化阶段，不能把“主线纳入”直接等同于“所有高级特性已经稳定落地”\citep{sglang_ascend_roadmap_2026}。

如果进一步参考这类项目的源码架构分析，可以把“主链复用 + 外围治理增强”的混合式路线抽象为更一般的组件分层：入口与启动护栏负责 CLI/OpenAI 服务和运行时预检，配置装配层负责把参数收敛为统一配置对象，引擎与 EngineCore 负责请求编排和调度通信，模型执行层负责 registry、worker、runner 与热点算子，平台与插件层负责能力探测和后端激活，多模态/Reasoning/Tool 层负责复杂能力的横切扩展，编译、KV Cache、分布式与 tracing 则构成性能基础设施\citep{vllmhust}。把这一结构放进综述并非为了做逐文件导读，而是因为它能帮助澄清混合式架构真正的工程含义：这一路线不是笼统地同时追求兼容与优化，而是通过分层把不同类型的国产化改动压回不同边界，使执行热路径优化、平台接入、环境治理和复杂能力扩展不必全部挤在同一条共享链路上。

从国产算力优化视角看，这种组件分层还有一个额外价值，即它能帮助识别“哪些组件最值得优先做本土化优化”。相比只给出系统级性能结论，表\ref{tab:hybrid-components} 更细地展示了混合式开源引擎各组件与国产算力潜在优化点之间的对应关系。这个视角对综述有意义，因为它把“混合式架构”的讨论从路线口号推进到了可操作的工程层：哪些部分更适合做 graph/compile 优化，哪些部分更适合做插件化平台适配，哪些部分更适合承载多模态与工具调用状态机，哪些部分则应该承担运行时护栏与环境治理。vLLM-HUST 在这里更适合被理解为这种组织方式的一个实例，而不是本节唯一需要关注的对象。

\begin{table*}[t]
\centering
\footnotesize
\setlength{\tabcolsep}{4pt}
\caption{混合式开源推理引擎组件分层与国产算力潜在优化点}
\label{tab:hybrid-components}
\begin{tabularx}{\textwidth}{>{\raggedright\arraybackslash}p{2.3cm}p{4.0cm}X}
\toprule
组件层 & 主要职责 & 可结合国产算力的潜在优化方向 \\
\midrule
入口与启动护栏 & CLI、OpenAI 兼容服务、启动时序控制、运行时预检与环境初始化 & 针对国产驱动/编译链构建更稳定的启动护栏，在引擎启动前完成 NPU 预检、环境变量自动补齐、依赖矩阵校验与故障前移，减少平台问题侵入执行主路径 \\
配置与能力装配 & 参数收敛、统一配置对象组装、不同能力开关与默认值选择 & 为国产后端提供平台感知的默认配置策略，如图模式开关、bucket 粒度、KV 容量预算、低比特路径和分布式后端选择，使优化从零散参数经验转为可复用配置模板 \\
引擎编排与 EngineCore & 请求对象转换、流式生命周期管理、调度与 core 通信 & 面向 shape 敏感后端优化连续批处理、prefill/decode 解耦、SLO 感知调度与前缀局部性管理，并在多进程/多实例形态下减少编译抖动和跨阶段干扰 \\
模型执行与热点算子 & 模型注册、权重加载、worker/runner、attention 与量化等热点路径 & 围绕国产硬件补齐 fused attention、MLA/MoE 相关算子、低比特 kernel、图安全回退路径和异步搬运策略，把平台特征转化为稳定吞吐和时延收益 \\
平台插件与运行时治理 & 平台探测、插件注入、后端激活、环境修复与外围治理工具衔接 & 把国产平台差异收敛为能力契约，通过插件化方式承载 dtype、attention backend、compile backend、distributed backend 与 runtime 修复逻辑，降低对共享主干的侵入深度 \\
多模态、Reasoning 与 Tool 横切层 & 多模态注册、IO 处理、reasoning parser、tool parser、复杂输出协议适配 & 把视觉 encoder、语音前端、结构化输出与工具调用状态统一纳入生命周期管理，并结合国产后端对 encode/decode 分段执行、跨模态缓存和 parser 前后处理开销做协同优化 \\
编译、KV Cache、分布式与 tracing 基础设施 & graph/compile、缓存管理、并行执行、监控与链路追踪 & 针对国产编译栈强化 graph cache 命中、piecewise backend 选择、KV 驻留层次、通信路径调优与近真实负载观测，使性能优化、容量管理和运行时诊断形成闭环 \\
\bottomrule
\end{tabularx}
\end{table*}

这一类生态已经不再只有少数头部系统，但真正值得关注的不是名单继续变长，而是开源 serving 生态正在沿几条相对清晰的机制路径分化。LightLLM、Text Generation Inference 与 MLC-LLM 更强调面向在线服务的高吞吐批处理、跨设备部署和统一服务接口，代表的是“优先优化服务主路径”的路线；Ollama、llama.cpp、ExLlamaV2 与 mistral.rs 更突出轻量化部署、本地运行和低比特执行，代表的是“优先保障受限资源下部署稳定性”的路线；Aphrodite、KTransformers 与 Lorax 则更关注运行时扩展、多租户适配和可定制 serving，代表的是“把推理引擎构造成可工程化组合平台”的路线\citep{lightllm,tgi,mlcllm,ollama,llamacpp,exllamav2,mistralrs,aphrodite,ktransformers,lorax}。这种分化对国产算力尤其重要，因为它提示我们：所谓“兼容开源生态”，并不是简单接入某个框架，而是要明确自己究竟优先继承哪一类能力主线，是吞吐优先、资源受限部署优先，还是多租户和平台化能力优先。

进一步看，开源主干框架在不同硬件生态上的演进方式也并不相同。NVIDIA 侧往往更快沉淀出围绕 Hopper、TensorRT-LLM、FlashAttention-3、FlexAttention 和 FlashInfer 的 kernel、编译优化和多模态运行时组合，其特点是把硬件代际能力迅速转化为 fused kernel 与 runtime 特性；AMD 侧当前更常见的则是围绕 MI300X、ROCm、vLLM 以及 SPX/DPX 分区模式组织多实例 serving、prefix caching 和前缀/解码阶段优化，其重点更偏向在现有框架边界内重组部署形态和资源切分\citep{tensorrtllm,flashinfer,shah2024flashattention3,dong2025flexattention,zhang2024vllmamd,song2026amdgpupartition,ambati2025mi300x}。它们虽然未必都直接面向国产算力，但共同构成了当前推理引擎生态的参照系，使国产后端适配必须同时回答“如何兼容主流接口”和“如何把本土平台特征转化为具体机制收益”这两个问题，而不只是把系统名字并列在一起。

\subsection{面向产业部署的一体化方案}

产业部署常采用“推理引擎 + 编译器 + 运行时 + 集群管理”一体化方案，以获得更强的可控性与稳定性。这类方案通常与特定芯片、编译链和部署平台深度绑定，能够对图编译、算子融合、通信调度和资源管理做更强约束，因此在某些固定模型与固定业务场景下往往具备更好的性能上限。但这种路线同样存在明显代价。首先，一体化系统往往需要为不同模型族、不同生成模式和不同上下文长度维护多组优化配置，其工程复杂度并不低。其次，这类平台若要兼容社区新模型或新服务特性，如工具调用、结构化输出和多 LoRA，往往需要较长的产品交付周期。换言之，一体化方案更适合追求交付稳定性与平台统一性的组织，而开源后端适配路线更适合希望快速跟进模型生态和上游能力迭代的团队。

在国产算力语境下，这一路线通常与本地基础软件栈深度耦合。MindSpore、CANN 与 MindIE 共同构成了较典型的国产软硬一体化推理路径，DeepSeek-V2 与 DeepSeek-V3 技术报告披露的 MLA、MoE 和大规模服务经验也表明，后端通信、缓存组织、注意力内核与平台运行时的一体化设计往往直接决定了系统可扩展上限\citep{mindspore,ascendcann,mindie,liu2024deepseekv2,deepseekv3}。与此相呼应，FlashMLA 展示了模型架构和底层注意力实现一体优化的工程方向，Mooncake 及其面向 KVCache 的存储解耦实践，则进一步说明产业界正在把缓存系统本身视作独立基础设施进行建设\citep{deepseekflashmla,qin2024mooncake,miao2024mooncake}。从横向比较看，无论是 NVIDIA 上围绕 Hopper 代际能力展开的 attention/kernel 深度优化，还是 AMD 围绕 MI300X 分区与 ROCm/vLLM 的部署工程实践，都说明一体化方案的性能上限往往来自对“芯片特征 + 运行时 + 服务策略”的联动打磨，而不是单点替换某个 kernel\citep{shah2024flashattention3,zhang2024vllmamd,song2026amdgpupartition,ambati2025mi300x}。

这也解释了为什么一体化方案在产业部署中往往具有更高的组织门槛。它要求团队不仅理解模型结构和推理算法，还要同时掌握编译器约束、运行时资源管理、设备通信语义以及部署平台的交付流程。对外表现上，这类系统常常提供的是一个“稳定可交付的平台能力包”，而非一个单独的 serving 框架；对内则意味着模型上线、版本回滚、性能调优与故障定位都被纳入同一套平台流程中。也正因如此，一体化路线的优势通常体现在平台统一性、固定场景确定性和长期运维可控性上，而不是对所有新模型和新特性的最快响应。从产业决策角度看，一体化方案的真正吸引力也并不只是跑分更高，而是能把很多本应由业务团队承担的系统复杂度前移到平台方内部。换言之，这类方案提供的不仅是推理性能，还包括版本矩阵收敛、交付责任集中和问题定位路径清晰等“组织确定性”。这也是为什么一体化方案在政企和行业部署中经常具有较强吸引力：在这些场景里，系统是否能被稳定运维、统一回归和持续交付，往往比单轮模型跟进速度更关键。相应地，它的代价也不只体现在生态兼容慢，更体现在系统创新往往需要通过平台产品节奏释放，外部团队很难像使用开源主干那样快速做局部替换和组合式实验。

与上述平台一体化方案相邻但不完全相同的，是 Chitu 与 xLLM 代表的国产原生/独立引擎路线。Chitu 的公开材料显示，其重点不是给 vLLM 增加一个后端，而是在独立生产级推理引擎内同时组织多元国产算力适配、FP8/FP4 低精度路径、专家并行（EP）、Prefill/Decode 分离（PD disaggregation）、Chunked Prefill、CPU+GPU 异构混合推理与集群部署能力\citep{chitu_tsinghua_news,chitu_github}。xLLM 则更突出 service-engine decoupled 架构：服务层承担在线/离线统一调度、动态 PD、Encode/Prefill/Decode 分离（EPD disaggregation）、多级 KV Cache 与容错治理，引擎层承担多流并行、图优化、显存管理、投机解码和 MoE 负载均衡；其公开硬件矩阵覆盖 NPU、MLU、ILU 与 MUSA 等多类国产或国产生态加速器，但模型覆盖和各硬件 feature 完整度并不均衡\citep{xllm_docs,xllm_service,xllm_tech_report}。这一路线对综述的价值在于说明，国产推理系统并非只能沿“主流框架加后端”演进，也可以从服务状态机、缓存治理、版本矩阵和部署闭环出发，自上而下设计一套独立系统边界。

\subsection{本土生态的关键分化点}

国产算力场景下，推理引擎的差异不再仅体现在单点算子性能，还体现在是否支持长上下文、Prefix Cache、Chunked Prefill、多 LoRA、多模态编码链路、结构化输出和工具调用等高层能力，以及这些能力能否在实际服务中稳定运行。特别是在真实服务中，单项能力“可用”与该能力“可规模化、可观测、可维护”之间存在显著差异。综合来看，本土生态中的技术分化主要体现在四个维度：其一，系统是优先追求通用框架兼容还是优先追求特定平台上的极致性能；其二，系统能否以较低代价跟进上游模型和推理特性的演进；其三，系统是否具备面向生产环境的灰度升级、故障回退与指标观测能力；其四，系统是否能够支持多样化工作负载，而不仅是固定 benchmark。这四个维度之所以重要，是因为它们对应了国产推理引擎最核心的四类成本：生态迁移成本、平台适配成本、交付运维成本和复杂负载成本。很多路线在单一维度上都能取得较好结果，但真正困难的是在多个维度上同时不过度失衡。例如，极端强调生态兼容的方案往往需要在平台深度优化上做妥协，极端强调单平台性能的方案又可能显著抬高后续模型迁移成本；而那些能力边界持续扩张的系统，若不能把这些支持组织成稳定抽象，最终仍会把成本转移为升级困难和维护负担。也正因此，国产推理引擎的比较不能只看功能清单，而要看一条路线究竟把成本压在了哪里、又把复杂度转嫁给了谁。

如果从当前已经能够坐实的公开主体和产品形态出发，本土路线至少可以粗分为三层：一类是沿开源主干或高性能工具链演进的后端适配与工具链强化路线，如 vLLM-Ascend、LMDeploy、飞桨系 FastDeploy，以及更强调高吞吐在线 serving 主路径优化的 LightLLM\citep{vllmascend,lmdeploy,fastdeploy,lightllm}；一类是把编译器、运行时、服务模板和交付流程整体收拢的平台一体化路线，如 MindIE；还有一类是直接以国产自研引擎或“自研框架 + 部署入口”形态组织能力边界的路线，如 Chitu 与 xLLM/玄武\citep{mindie,chitu_tsinghua_news,chitu_qingcheng,chitu_github,xllm_docs,xuanwu_tsingmao}。vLLM-HUST 则更接近混合式工程组织：它不完全等同于独立引擎，也不只是后端适配，而是在尽量复用上游服务主体的同时，把平台护栏、插件链和外围治理工具沉淀为独立工程层\citep{vllmhust}。这一分层之所以重要，不是为了扩充名单，而是为了避免把“框架本体”“交付平台”和“部署入口”混成同一层对象，从而误判各路线真正承担的复杂度类型。

还需要再向前分清一层：公开搜索里经常会同时出现工作站、一体机、云平台、训推平台和应用交付页面，但这些对象大多是在封装、承载或预集成上述引擎路线，而不是与推理引擎本体处于同一比较层级。换言之，部署载体回答的是“能力通过什么形态被交付出去”，推理引擎回答的则是“请求、状态、缓存、调度和算子如何被组织起来”。如果把两者直接并列，结论就会从路线比较滑向产品罗列，既不利于解释系统复杂度真正落在哪一层，也容易把工程入口便利性误写成引擎内核先进性。

\begin{table*}[t]
\centering
\footnotesize
\setlength{\tabcolsep}{4pt}
\caption{国内相关组织与国产算力推理引擎路线定位}
\label{tab:domestic-org-route}
\begin{tabularx}{\textwidth}{>{\raggedright\arraybackslash}p{2.8cm}p{2.5cm}p{3.5cm}X}
\toprule
公开主体或组织生态 & 代表对象 & 更接近的技术控制点 & 对综述主题的启示 \\
\midrule
清华大学相关团队与清程极智 & Chitu「赤兔」\citep{chitu_tsinghua_news,chitu_qingcheng,chitu_github} & 独立生产级引擎、多元算力适配、PD/EP 与低精度路径 & 高校与产业团队可以直接经营自研执行栈，把多硬件适配和生产部署能力放进同一引擎边界 \\
京东开源团队 & xLLM / xLLM-Service\citep{xllm_docs,xllm_service,xllm_tech_report} & 服务层--引擎层分离、动态 PD/EPD、多级 KV Cache 与状态治理 & 大厂开源路线更容易从服务治理、缓存治理和部署闭环反向塑造引擎结构 \\
华为昇腾生态与开源社区 & MindIE、vLLM-Ascend、SGLang Ascend\citep{mindie,vllmascend,sglang_ascend_npu} & 平台一体化交付与主流 serving 框架后端吸纳并行 & 国产 NPU 既可以通过平台产品栈交付，也可以通过 plugin backend 或主线平台后端进入国际开源生态 \\
寒武纪生态 & vLLM-MLU\citep{vllmmlu} & 基于 vLLM plugin system 的 MLU 后端适配 & 芯片厂商路线的关键不只在算子补齐，还在 SDK/runtime、分布式组件和上游版本对齐 \\
飞桨 / PaddlePaddle 生态 & FastDeploy / PaddleNLP Server\citep{fastdeploy,fastdeploy_docs,paddlenlp_llm_server} & 多硬件部署工具链、服务接口、KV 传输与量化部署 & 深度学习框架生态可以把国产硬件适配沉淀为“工具链 + 服务封装 + 安装入口”的工程体系 \\
OpenMMLab / InternLM 生态 & LMDeploy\citep{lmdeploy,lmdeploy_install,lmdeploy_supported_models} & TurboMind 与 PyTorchEngine 双后端工具链 & 同一部署工具包在 CUDA 与国产平台上可能走不同后端路径，性能结论和功能矩阵不应跨平台外推 \\
ModelTC / LightLLM 团队 & LightLLM\citep{lightllm,lightllm_docs_2026,gong2025past} & 高吞吐 serving、TokenAttention、调度与内核生态 & 国内团队也在贡献通用 serving 机制，但公开证据更支持其作为生态补充而非多国产硬件主线引擎 \\
中国电信研究院及合作方 & Triton 统一跨架构推理框架\citep{chinatelecom_triton_validation_2025,chinatelecom_wbba_2026} & 跨架构编译器、统一算子库、透明运行时插件 & 运营商/平台方关注多芯重复适配与统一迁移，但未开源和复现边界决定其更适合作为跨架构抽象探索 \\
\bottomrule
\end{tabularx}
\end{table*}

表\ref{tab:domestic-org-route} 的意义不在于给出一个“单位名录”，而在于揭示组织位置会直接影响技术路线：芯片与平台生态更倾向把复杂度收进 SDK、runtime、版本矩阵和平台交付链路；框架与工具链团队更倾向把复杂度压到 plugin backend、服务接口和模型支持矩阵；运营商或大型平台方则更容易从异构资源池运营出发，推动跨架构编译和运行时抽象。由此可见，面向国产算力的推理引擎竞争并不是单一项目之间的横向竞赛，而是“谁掌握抽象边界、谁承担适配成本、谁负责交付闭环”的组织分工问题。

在这一分层之上，还应把“组织正在研究或交付什么”与“哪些能力已经能被公开证据坐实”分开。可公开核验的材料大致形成了四类证据强度：源码、安装文档和版本矩阵较清楚的 plugin backend 或工具链项目，可以支撑“国产平台正在进入主流 serving 接口与部署流程”的判断；技术报告、服务层仓库和官方文档较充分的 Chitu 与 xLLM，可以支撑“独立引擎正在从状态机、KV Cache 与部署闭环重构国产推理系统边界”的判断；MindIE 与中国电信 Triton 这类平台或跨架构方案，可以支撑“平台方尝试把复杂度收进 SDK/runtime 或编译运行时抽象”的判断，但不宜写成完全开源、可复现实验体系；LightLLM 则更适合证明国内团队参与高吞吐 serving 机制创新，而不是证明其已经成为国产多硬件主线引擎\citep{vllmmlu,sglang_ascend_npu,fastdeploy_docs,lmdeploy_install,chitu_github,xllm_tech_report,mindie,chinatelecom_triton_validation_2025,lightllm_docs_2026}。因此，本综述后续比较应同时标注三条边界：能力是否进入主线或公开仓库，版本与硬件矩阵是否可核验，性能或 roadmap 表述是否来自可复现实验。这个边界意识比继续增加项目数量更能提升本节的审稿可信度。

公开社区讨论虽然不能替代系统论文、官方文档和源码分析，但对理解本土路线为何分化仍有补充价值，因为它更直接暴露了团队在真实部署中最常遇到的摩擦点。与此同时，清华大学与清程极智联合开源的赤兔 Chitu 说明，本土生态并不只是在既有主干和平台栈之间做二选一，也在尝试以独立生产级引擎直接组织多元算力兼容、弹性部署和企业落地能力\citep{chitu_tsinghua_news,chitu_qingcheng,chitu_github}。xLLM 官方文档则把自身明确定位为面向国产芯片的开源智能推理框架，而玄武 XuanWu 页面进一步强调一键部署、OpenAI 兼容和推理引擎预验证等工程入口，说明国内也在形成“推理框架 + 部署入口”一体化演进的产品形态\citep{xllm_docs,xuanwu_tsingmao}。与此同时，LightLLM 这类高吞吐 serving 项目的存在也提示我们，国内团队并不只是在“国产后端适配”或“平台一体化交付”两个方向上展开工作，还在直接参与在线服务主路径、批处理组织和吞吐导向运行时的机制演进\citep{lightllm}；而 FastDeploy 与 PaddleNLP LLM Server 这类飞桨系入口，则更清楚地展示了另一种“工具链强化 + 服务封装”路线，即不把重点放在重写通用 serving 主体，而是把模型导出、高性能推理与服务接入组织成更完整的部署链路\citep{fastdeploy,paddlenlp_llm_server}。2025--2026 年的知乎部署实战帖反复强调三件事：其一，vLLM-Ascend 的工程可行性高度依赖“上游主干版本 + 插件分支 + CANN/torch\_npu”严格对齐，且很多经验都围绕如何在不改上层调用代码的前提下借助插件化后端完成迁移\citep{zhihu_vllm_ascend_guide,zhihu_vllm_pd_migration}；其二，MindIE 的吸引力并不只来自单点推理性能，而在于它把镜像、服务化组件、监控和配置模板一起交付，从而把一部分系统复杂度前移到平台栈内部\citep{zhihu_mindie_deploy}；其三，PD 分离、多机 RoCE/HCCL 连通性、环境预检和日志/配置模板在社区记录中的出现频率极高，说明国产平台的一线痛点往往是“先把系统稳定组织起来”，之后才谈局部极致优化\citep{zhihu_vllm_pd_migration}。

工具链型路线进一步说明，国产算力适配不只发生在单一 runtime 的后端层。FastDeploy 更接近飞桨生态中的多硬件推理部署工具链：它把 OpenAI/vLLM 风格服务接口、PD 分离、KV Cache 传输、量化、投机解码、Chunked Prefill 与按硬件划分的安装入口组织在同一工程体系中，面向昆仑芯、海光、昇腾、天数、燧原、沐曦等平台提供公开适配入口\citep{fastdeploy,fastdeploy_docs}。LMDeploy 则需要明确区分 TurboMind 与 PyTorchEngine：前者主要服务 CUDA 高性能路径，国产平台扩展主要通过 PyTorchEngine/Other Platforms 文档进入 Ascend、Cambricon 和 MACA 等后端，并且模型、设备与 eager/graph/量化模式的支持粒度并不对称\citep{lmdeploy,lmdeploy_install,lmdeploy_supported_models}。因此，二者都不宜被简单写成“又一个推理引擎”，而更适合作为“部署工具链 + 服务接口 + 硬件适配层”的工程体系样本。

最后还应保留一个次级但有趋势意义的生态位置：LightLLM 更适合作为国内团队在高吞吐 serving、TokenAttention、Efficient Router、调度与内核生态上的重要补充，而不是“成熟国产多硬件主线引擎”；其公开文档中最明确的国产硬件线索主要是 MUSA 路径，默认主线仍明显偏 CUDA/Triton\citep{lightllm,lightllm_docs_2026,gong2025past}。与此不同，中国电信 Triton 统一跨架构推理框架并不处于 vLLM、LMDeploy 或 Chitu 同一层，而更像跨架构编译/运行时抽象：官方材料强调自研 Triton 跨架构编译器、统一算子库、vLLM-Triton 透明嵌入插件和图算融合编译器，用于缓解多国产芯片重复适配问题；但由于代码、许可证、版本矩阵和复现脚本尚未公开，相关迁移时间与性能损耗数字只能作为官方技术验证口径谨慎引用\citep{chinatelecom_triton_validation_2025,chinatelecom_wbba_2026}。

这一观察与公开社区讨论中的高频主题基本一致。围绕“Ascend vllm”和“国产算力 推理引擎”，反复出现的是原生支持、双机组网、推理优化、MindIE 入门和国产化实践等议题。尽管这些讨论未必提供可复现实验，但足以揭示社区的实际判断：路线比较首先看环境护栏、网络链路、配置模板、服务化接口和回退路径是否完备，而不是单卡跑分是否领先。对综述而言，这类“社区摩擦点”补上了论文和官方材料不易充分覆盖的工程现实，即一条路线能否持续，往往先取决于其是否降低部署不确定性。

将这些社区信号收敛为选型判断，比直接比较“谁更快”更有实用性：vLLM-Ascend 适合已有 vLLM 服务栈、希望少改上层接口并快速验证国产迁移的团队，但需要接受严格版本矩阵与插件同步节奏；MindIE 适合把统一交付、模板化配置、镜像与平台监控放在首位的组织，其代价是外部可塑性较低；Chitu 更适合希望直接围绕独立生产级引擎组织多元算力兼容与企业落地能力的团队；xLLM / 玄武更接近“推理框架 + 部署入口”一体化思路，适合希望把国产芯片适配经验进一步沉淀为统一操作入口和部署闭环的场景；LMDeploy 适合兼顾模型覆盖、部署效率和高性能工具链的场景，但国产 GPU 适配成熟度仍需逐平台核验；SGLang 在中文社区更常对应复杂 Agent、长 system prompt 复用和结构化输出场景，其优势主要来自对复杂执行状态的运行时组织；vLLM-HUST 类混合式路线则适合既要维持上游接口兼容、又要逐步沉淀国产平台护栏与治理能力的团队\citep{zhihu_vllm_ascend_guide,zhihu_vllm_pd_migration,zhihu_mindie_deploy,zhihu_lmdeploy_domestic,zhihu_vllm_sglang_llamacpp_compare,xhs_search_lmdeploy,xhs_search_mindie,xhs_search_sglang_deploy,chitu_tsinghua_news,chitu_qingcheng,chitu_github,xllm_docs,xuanwu_tsingmao}。

\begin{table*}[t]
\centering
\footnotesize
\setlength{\tabcolsep}{4pt}
\caption{结合社区实战信号的国产推理引擎路线选型提示}
\label{tab:community-route-selection}
\begin{tabularx}{\textwidth}{>{\raggedright\arraybackslash}p{2.0cm}p{2.2cm}p{3.6cm}p{3.6cm}X}
\toprule
代表对象 & 更适合的前提 & 社区中反复暴露的工程摩擦 & 更值得关注的系统收益 & 更接近的组织选择 \\
\midrule
vLLM-Ascend & 已有 vLLM 或 OpenAI 兼容服务栈，希望优先继承上游模型支持与接口形态 & 上游主干、插件分支、CANN 与 torch\_npu 需要严格对齐；PD 分离、多机组网和环境预检经常成为首轮障碍\citep{zhihu_vllm_ascend_guide,zhihu_vllm_pd_migration,xhs_search_vllm_ascend} & 以较低上层改动成本验证国产迁移路径，较快吸收上游 serving 特性 & 生态跟进速度优先、能承受版本协同与环境治理成本的团队 \\
MindIE & 更重视镜像、模板、监控、服务组件和平台统一交付，希望把部署复杂度收敛到平台栈内部 & 配置参数、服务模板、平台组件依赖和网络环境准备对部署成功率影响很大\citep{zhihu_mindie_deploy,xhs_search_mindie} & 交付责任集中、运维路径清晰、平台确定性更强 & 行业部署、政企交付或平台统一性优先的组织 \\
Chitu & 希望直接采用独立生产级引擎组织多元算力兼容、弹性部署和企业落地能力\citep{chitu_tsinghua_news,chitu_qingcheng,chitu_github} & 独立引擎需要同时经营模型支持矩阵、服务接口和交付模板，若模块边界不清晰，演进成本会快速上升 & 更容易把多硬件兼容、生产可用性和部署弹性放到同一套自研内核里统筹 & 既希望摆脱单一上游主干依赖、又具备持续经营自研内核能力的团队 \\
xLLM / 玄武 & 希望围绕国产芯片同步建设推理框架与部署入口，把一键拉起、OpenAI 兼容与预验证能力前置到工程闭环中\citep{xllm_docs,xuanwu_tsingmao} & 需要同时维护推理内核、分布式执行与部署入口的一致性，避免把入口便利性建立在脆弱的环境假设上 & 更容易形成面向国产芯片的“框架 + 部署入口”一体化落地路径 & 希望把国产化部署经验沉淀为统一操作入口和服务闭环的团队 \\
LMDeploy & 同时看重推理性能、模型支持矩阵和工程化工具链，希望保留较强部署灵活性 & 国产 GPU 适配成熟度、量化路径、不同硬件上的最佳实践需要逐平台核验，社区比较常把它放在与 vLLM 的横评语境中\citep{zhihu_lmdeploy_domestic,xhs_search_lmdeploy} & 部署效率和高性能工具链较平衡，适合快速工程试验与模型接入 & 追求性能与工程效率折中、愿意做平台侧验证的团队 \\
SGLang & 更重视复杂 Agent、长上下文前缀复用、结构化输出和请求编排，希望把复杂控制流直接写进 serving 运行时 & 中文社区更常把它放在与 vLLM 的取舍语境中，但国产化部署模板和硬件侧现成经验相对少，工程门槛更多落在负载理解与执行模式选择上\citep{zhihu_vllm_sglang_llamacpp_compare,xhs_search_sglang_deploy} & 在复杂控制流、长 system prompt 复用和结构化输出场景下更容易体现运行时组织优势 & 复杂工作负载优先、愿意为执行模型与调度抽象投入理解成本的团队 \\
vLLM-HUST 类混合式路线 & 不愿放弃上游接口兼容，同时希望把国产平台护栏、插件与治理工具逐步沉淀为独立工程层 & 需要持续经营模块边界，避免既继承上游复杂度又堆积本地特化补丁；对工程架构能力要求更高\citep{vllmhust} & 在接口兼容、平台特化和运行时治理之间形成可持续折中 & 既要长期跟进上游生态，又要持续建设本地化治理能力的团队 \\
\bottomrule
\end{tabularx}
\end{table*}

从综述写作角度看，表\ref{tab:community-route-selection} 更适合回答“什么样的团队在什么前提下会倾向哪条路线”；而第 4 节的表\ref{tab:route-challenge-mapping} 则进一步回答“选了这条路线之后，首先会在哪些系统问题上付出代价”。两表合读时，本节讨论就不再是松散的框架介绍，而是在说明国产推理引擎路线选择本质上是一种复杂度分配：有的路线把复杂度压在版本矩阵和平台适配上，有的路线压在平台交付与模板治理上，有的路线压在复杂控制流和运行时状态空间上，还有的路线压在模块边界与长期 merge-safe 演进上。

围绕多租户与定制化模型服务的能力也在持续拉开系统差距。Punica 与 S-LoRA 表明，当 LoRA 适配器数量迅速增长时，系统需要在缓存复用、批内共享和运行时隔离之间重新平衡\citep{chen2023punica,slora2024}。与此同时，Guidance、Outlines 与 XGrammar 等工具说明，结构化输出与受约束生成正在把“推理引擎”从纯粹的 token 生成器转变为可编排的执行系统；而 XGrammar 的论文实现则进一步表明，这种能力已经需要以语法执行引擎、持久化栈和 GPU 重叠执行等系统机制来支撑\citep{guidance,outlines,xgrammar,dong2024xgrammar}。

调度能力同样正在成为新的分化点，但这里更重要的也不是系统名本身，而是它们分别把哪一类调度变量显式化了。Llumnix 把跨实例动态迁移和内存状态搬运带入运行时，使“负载均衡”不再只是请求排队问题；SOLA 把 TTFT、TPOT 等服务级目标直接写进调度目标函数，使“调度最优”从吞吐优先转向 SLO 优先；AI Metropolis 则进一步把多 Agent 依赖关系显式化，通过乱序执行削减伪串行等待\citep{sun2024llumnix,hong2025sola,xie2025aimetropolis}。这一组工作共同说明，现代推理系统的竞争焦点已经从“是否支持连续批处理”进一步演进到“是否能够在复杂状态空间里稳定组织执行”，而这正是国产推理引擎在复杂业务负载下最不能回避的技术点。

\begin{table}[t]
\centering
\small
\caption{国产算力推理引擎技术路线比较}
\begin{tabularx}{\linewidth}{>{\raggedright\arraybackslash}p{2.2cm}YYY}
\toprule
路线 & 代表思路 & 主要优势 & 主要代价 \\
\midrule
开源主干后端适配 & 基于 vLLM、SGLang 等主干框架增加国产平台后端 & 复用上游生态快，模型覆盖广，便于跟踪新特性 & 需要控制 patch 深度，维护抽象边界 \\
平台一体化引擎 & 推理引擎、编译器、运行时与部署平台协同优化 & 平台可控性强，固定场景性能上限高 & 生态兼容慢，长期维护成本高 \\
混合式架构 & 上层接口兼容开源生态，底层热点路径深度定制 & 在兼容性与性能之间取得折中 & 体系复杂，对工程团队要求高 \\
\bottomrule
\end{tabularx}
\end{table}

图\ref{fig:route-divergence} 将三类路线放到同一张图里比较，其重点不在功能罗列，而在它们对“生态兼容性、平台特化深度、交付稳定性”三者的不同取舍\citep{vllmhust}。从课题化建设视角看，这种分化本质上也是对执行与通信、状态治理和压缩协同三条主路径控制权的不同分配，因此更能解释为什么一些系统在 benchmark 上接近，却在长期维护和场景扩展上差异巨大。

\begin{figure}[t]
\centering
\resizebox{0.98\linewidth}{!}{\begin{tikzpicture}[
    node distance=5mm and 10mm,
    block/.style={draw, rounded corners=2pt, minimum width=31mm, minimum height=9mm, align=center, font=\small, inner sep=4pt},
    note/.style={draw, rounded corners=2pt, minimum width=31mm, align=center, font=\scriptsize, inner sep=3pt},
    flow/.style={-{Latex[length=2.2mm]}, semithick}
]
\node[block, fill=gray!10, minimum width=42mm] (goal) {国产推理引擎路线选择};
\node[font=\scriptsize, align=center, above=1mm of goal] {核心张力：生态兼容性、平台特化深度、交付稳定性};

\node[block, fill=blue!10, below left=12mm and 31mm of goal] (upstream) {开源主干后端适配\\vLLM / SGLang / vLLM-Ascend};
\node[block, fill=orange!12, below=12mm of goal] (hybrid) {混合式架构\\vLLM-HUST\\上层兼容 + 底层热点特化};
\node[block, fill=green!10, below right=12mm and 31mm of goal] (integrated) {一体化方案\\MindIE / 编译器 / 运行时协同};

\node[note, fill=blue!5, below=of upstream] (u1) {优势：生态吸收快\\模型跟进快};
\node[note, fill=red!5, below=of u1] (u2) {代价：patch 累积\\抽象边界易失稳};

\node[note, fill=orange!6, below=of hybrid] (h1) {优势：兼顾兼容性\\与平台优化};
\node[note, fill=red!5, below=of h1] (h2) {代价：体系复杂\\组织门槛高};

\node[note, fill=green!5, below=of integrated] (i1) {优势：交付稳定\\平台约束强};
\node[note, fill=red!5, below=of i1] (i2) {代价：生态兼容慢\\上游响应较慢};

\draw[flow] (goal) -- (upstream);
\draw[flow] (goal) -- (hybrid);
\draw[flow] (goal) -- (integrated);
\draw[flow] (upstream) -- (u1);
\draw[flow] (u1) -- (u2);
\draw[flow] (hybrid) -- (h1);
\draw[flow] (h1) -- (h2);
\draw[flow] (integrated) -- (i1);
\draw[flow] (i1) -- (i2);
\end{tikzpicture}}
\caption{国产推理引擎技术路线分化示意图}
\label{fig:route-divergence}
\end{figure}

从现实可行性看，混合式路线更符合当前阶段，因为它试图同时保留上游生态复用与本土关键路径特化。但它成立的前提，是上层生态节奏与底层平台节奏之间存在清晰边界；否则系统很容易演化为既继承上游复杂度、又承受本地特化负担的双重分叉。

\begin{table}[t]
\centering
\footnotesize
\setlength{\tabcolsep}{4pt}
\caption{典型推理引擎能力比较}
\begin{tabularx}{\linewidth}{>{\raggedright\arraybackslash}p{2.1cm}p{1.8cm}p{1.9cm}p{1.8cm}Y}
\toprule
系统或路线 & 长上下文/缓存 & 多模态/结构化输出 & 平台绑定程度 & 更适合的使用情形 \\
\midrule
vLLM 类主干框架 & 强，强调 PagedAttention 与连续批处理\citep{vllm,kwon2023efficient} & 持续增强，迭代快 & 低到中 & 希望快速跟进上游模型与服务特性的团队 \\
vLLM-HUST & 强，继承 vLLM 主干并强化插件化平台接入\citep{vllmhust} & 强，覆盖 OpenAI 接口、多模态与 Multi-LoRA 等工程能力 & 中 & 既要跟进上游生态，又要补齐国产平台接入与运行时治理的团队 \\
SGLang 类执行框架 & 强，强调复杂请求编排与服务吞吐\citep{sglang} & 强，适合 Agent 与复杂控制流 & 低到中 & 重视复杂负载调度与执行效率的团队 \\
MindIE 类一体化方案 & 中到强，依赖平台优化深度\citep{mindie} & 中，能力受产品路线影响 & 高 & 平台统一、交付稳定优先的行业部署 \\
Chitu「赤兔」 & 强，强调多元算力适配与生产级部署\citep{chitu_qingcheng,chitu_github} & 中到强，取决于独立框架能力边界 & 中 & 希望在多硬件环境下维持统一推理框架并平滑扩展部署规模的团队 \\
xLLM / 玄武类路线 & 强，强调国产芯片适配、分布式推理与部署入口一体化\citep{xllm_docs,xuanwu_tsingmao} & 中，当前更突出推理与部署闭环 & 中到高 & 希望围绕国产芯片形成自研框架与部署闭环的团队 \\
LMDeploy 等高性能工具链 & 强，热点路径优化充分\citep{lmdeploy,tensorrtllm} & 中到强，随版本持续增强 & 中到高 & 追求高性能部署与工程化工具链的场景 \\
LightLLM 类高吞吐 serving 路线 & 强，强调在线服务吞吐、批处理组织与 serving 主路径优化\citep{lightllm} & 中，当前公开材料更突出文本服务主路径 & 低到中 & 希望围绕高吞吐在线 serving 持续打磨执行主路径的团队 \\
\bottomrule
\end{tabularx}
\end{table}

需要指出的是，上述能力对比并非静态结论，而是会随着上游版本迭代、国产后端接入深度和具体部署目标持续变化。综述层面更值得关注的是一种判断框架：如果目标是快速吸收社区能力，应优先选择抽象边界清晰、插件接口完善的路线；如果目标是固定平台上的确定性交付，则更应重视图编译、运行时和运维体系的一体化程度。

\begin{table}[t]
\centering
\scriptsize
\setlength{\tabcolsep}{4pt}
\caption{国内部署路线案例比较}
\begin{tabularx}{\linewidth}{>{\raggedright\arraybackslash}p{2.0cm}p{1.7cm}p{2.7cm}Y}
\toprule
代表对象 & 路线类型 & 工程落点 & 部署目标 \\
\midrule
vLLM-Ascend\citep{vllmascend} & 开源主干后端适配 & 围绕上游接口扩展后端、执行器与定制算子 & 快速跟进上游能力，优先保障生态兼容与模型覆盖 \\
vLLM-MLU\citep{vllmmlu} & 开源主干后端适配 & 依托 vLLM plugin system 接入寒武纪 MLU，并覆盖 Chunk Prefill、Prefix Caching、Spec Decode 等 vLLM 能力 & 以继承 vLLM 接口与调度语义为主，但复现性仍受 SDK/runtime 与公开 benchmark 粒度约束 \\
SGLang Ascend\citep{sglang_ascend_npu,sglang_ascend_roadmap_2026} & 主线后端支持 & 在 SGLang 主线中为 Ascend NPU 提供 attention backend、PD 分离部署和组件版本映射 & 体现国产 NPU 被主流 serving 框架吸纳，但高级能力成熟度需与 roadmap 区分 \\
vLLM-HUST\citep{vllmhust} & 混合式架构 & 保持 vLLM 上层接口，外挂平台插件并补齐运行时护栏与环境治理 & 兼顾上游兼容、国产平台接入与部署运维闭环 \\
MindIE\citep{mindie} & 一体化部署 & 与本地编译器、运行时和交付链路深度协同 & 平台统一交付、稳定运行和行业方案落地优先 \\
Chitu「赤兔」\citep{chitu_tsinghua_news,chitu_qingcheng,chitu_github} & 独立自研引擎 & 面向多元算力组织生产级推理、弹性部署与企业落地能力 & 同时重视多硬件兼容、部署弹性与生产可用性的场景 \\
xLLM / 玄武\citep{xllm_docs,xuanwu_tsingmao} & 自研框架 + 部署入口 & 围绕国产芯片组织开源智能推理框架，并以前端部署入口补齐一键拉起、OpenAI 兼容与预验证能力 & 希望同步建设推理内核与部署闭环的国产化落地场景 \\
LMDeploy\citep{lmdeploy} & 工具链强化 & 围绕模型支持、推理工具链和工程化部署体验持续扩展 & 兼顾高性能部署、模型覆盖和工程效率 \\
FastDeploy / PaddleNLP Server\citep{fastdeploy,paddlenlp_llm_server} & 工具链强化 + 服务封装 & 以高性能推理与部署工具包衔接模型导出、服务封装和应用接入 & 希望围绕飞桨生态组织 LLM/VLM 推理与部署链路的场景 \\
中国电信 Triton 统一跨架构框架\citep{chinatelecom_triton_validation_2025} & 跨架构编译/运行时 & 通过 Triton 跨架构编译器、统一算子库和 vLLM-Triton 透明插件探索多芯迁移 & 用于缓解多芯重复适配问题，但未开源且许可证与复现脚本未明确 \\
\bottomrule
\end{tabularx}
\end{table}

多模态推理正在把引擎边界从文本 decode 推向异构阶段编排。就国内团队的工程案例看，Qwen2-VL 与 InternVL2 代表了多图与文档理解链路进入在线服务后的典型压力点：系统不仅要处理动态分辨率切图和图像 token 打包，还要面对视觉前缀复用不足、OCR 密集区域导致的 prefill 膨胀，以及图文混排输入带来的 shape 波动\citep{qwen2vl,internvl,wang2024qwen2vl,chen2024internvl2}。CogVLM2 则把视频理解显式纳入主流模型服务，使推理链路必须同时考虑帧采样、跨帧压缩、视觉 encoder 批处理以及 decode 侧连续 batching 之间的资源再分配\citep{cogvlm2,wang2024cogvlm2}。在语音方向，Qwen2-Audio、GLM-4-Voice 与 MiniCPM-o 进一步把音频前端、speech tokenizer、流式 chunk 累积以及语音/文本交错生成纳入同一运行时，使首 token 时延、实时率和打断恢复直接成为推理引擎的系统约束\citep{chu2024qwen2audio,zeng2024glm4voice,cui2026minicpmo45}。

如果从运行时组织角度进一步拆解，这些压力并非简单叠加在文本 serving 主路径之上，而是在三个层次上重写引擎设计。第一，多图、文档与 OCR 负载首先重写的是 prefill 主路径：视觉 token 数量膨胀、切图策略离散化以及图文混排导致的长度抖动，会直接改变 bucket 划分、图模式命中率与前缀缓存复用策略。第二，视频负载重写的是阶段关系本身：帧采样、视觉 encoder 和语言 decode 不再是线性串接关系，而是必须围绕显存峰值、带宽占用和首 token 时延做跨阶段再平衡。第三，语音交互重写的是会话状态机：streaming chunk、实时打断、语音文本交错生成以及会话恢复，使系统不再只管理 KV Cache，还要同时管理音频前端状态、增量转写结果和流式输出节奏。也就是说，多模态路径真正扩张的不是“模型输入类型”，而是引擎内部必须长期维护的状态对象种类。

这也是为什么多模态引擎往往不能通过“在文本框架前面补一层预处理”来完成工程闭环。对于多图文档理解，系统需要把图像切块、版面解析和视觉前缀编码结果纳入可缓存对象，否则同一文档的重复查询难以获得稳定复用收益；对于视频理解，系统需要显式决定哪些阶段可以批内共享、哪些中间结果应跨帧复用、哪些请求必须在 encode 与 decode 之间做资源隔离；对于语音对话，则要把实时率、打断恢复和声文混合生成直接写进调度与回退逻辑。正因为这些约束同时涉及前处理、缓存、调度、图执行和服务协议，多模态推理才会成为检验一条引擎路线是否成熟的分水岭：如果系统只能做到单模型单链路跑通，而不能把这些异构阶段收敛为统一状态机，那么模型支持矩阵越扩张，运行时复杂度就越难被控制。

\begin{table}[t]
\centering
\footnotesize
\setlength{\tabcolsep}{4pt}
\caption{国内团队多模态推理链路工程关注点比较}
\begin{tabularx}{\linewidth}{>{\raggedright\arraybackslash}p{2.2cm}p{1.9cm}p{2.2cm}Y}
\toprule
代表模型 & 主要链路 & 关键工程压力 & 对推理引擎的直接要求 \\
\midrule
Qwen2-VL、InternVL2\citep{wang2024qwen2vl,chen2024internvl2} & 多图、文档与 OCR & 动态分辨率切图、视觉 token 膨胀、图文混排带来 shape 波动 & 视觉前处理与文本 prefill 协同、前缀缓存复用、bucket 化与编译抖动控制 \\
CogVLM2\citep{wang2024cogvlm2} & 视频理解 & 帧采样与跨帧压缩策略影响 encode/decode 资源分配 & 视觉 encoder 批处理、跨阶段资源再平衡、视频前缀组织 \\
Qwen2-Audio 等\citep{chu2024qwen2audio,zeng2024glm4voice,cui2026minicpmo45} & 语音与实时交互 & 音频前端、speech tokenizer、流式 chunk 与打断恢复共同影响时延 & 流式调度、语音文本交错生成、实时率控制与会话状态恢复 \\
\bottomrule
\end{tabularx}
\end{table}

ElasticMM、EPDServe 与 Eevee 分别从弹性并行、阶段解耦和模块复用三条路径回答这些问题：前者统一管理多模态请求与前缀缓存，后两者分别通过 encode/prefill/decode 拆分和 module multiplexing 缓解阶段异质性与模块争用\citep{liu2025elasticmm,singh2025epdserve,hong2026eevee}。它们共同说明，多模态系统的难点已不是“支持某个模型”，而是能否在请求异构、阶段异质和模块干扰并存时保持资源利用率与首 token 时延稳定。对国产推理引擎而言，这要求把前处理、跨模态缓存、阶段解耦和异构编码链路一起纳入运行时设计。

进一步看，这些工作之所以对国产路线具有参考价值，不是因为它们提供了可以直接照搬的实现，而是因为它们把多模态运行时中最容易被忽略的几个控制变量显式化了。ElasticMM 强调多模态请求不应只在入口处分类，而应在运行时持续维护视觉前缀、缓存命中和并行度之间的关系；EPDServe 说明 encode、prefill 与 decode 一旦继续共用同一套排队和资源分配策略，视频和文档类请求就会系统性挤压文本 decode 的服务质量；Eevee 则进一步表明，模块复用并不是单纯节省显存，而是在异构阶段交替出现时减少重复加载与执行抖动的一种运行时组织方法\citep{liu2025elasticmm,singh2025epdserve,hong2026eevee}。对国产推理引擎而言，这意味着多模态工程的关键不只是补齐若干模型适配器，而是把视觉 encoder、语音前端、结构化输出、tool/reasoning parser 与平台后端能力一起压进统一的生命周期管理中。只有当这些对象都能够被调度器、缓存管理器和异常回退链共同感知，本土系统才可能在多模态场景下同时维持吞吐、时延和可维护性。

以 vLLM-HUST 这类本土 fork 为例，其工程重点已经不再只是补一层国产平台接入，而是把多模态注册中心、reasoning parser、tool parser、OpenAI 兼容服务与平台插件链一起组织成横切能力层，使模型接入、工具调用和复杂输出协议能够在同一 serving 主体内持续演进\citep{vllmhust}。这类做法说明，国产推理引擎的“可持续性”越来越依赖能力边界是否模块化，而不只是某一条 kernel 路径是否被局部打磨到极致。

从本土生态的工程走向看，一个越来越清晰的趋势是：模型能力边界正在反向塑造推理引擎边界。无论是 LMDeploy 围绕部署工具链与模型支持矩阵的持续扩展，还是 MindIE 一类一体化系统对平台统一性交付的强调，抑或 vLLM-Ascend 这类后端适配项目对上游抽象边界的依赖，其共同约束都不再只是“文本生成是否足够快”，而是多模态、结构化输出、长上下文和复杂工作负载能否被纳入同一套可维护的运行时框架\citep{lmdeploy,mindie,vllmascend}。

换言之，当前国产推理引擎的比较维度已经从“谁支持更多模型”逐步转向“谁能以更低工程代价吸收模型变化”。当上游模型不断引入 MoE、长上下文、多模态 encoder 和语音交互等新结构时，真正稀缺的能力不再是一次性的端到端适配，而是能否把这些变化收敛成可复用的调度、缓存、算子与接口演进机制。对综述而言，这一点比单次性能对比更重要，因为它更直接决定了本土生态能否形成持续积累而非反复重写的工程路径。因此，本节真正想给出的并不是某一条固定答案，而是一组更细的判别准则：如果团队首要目标是快速跟进模型与服务特性，应优先选择上游距离近、插件接口清晰且共享路径修改受控的方案；如果目标是面向固定业务长期稳定交付，则应更重视一体化平台对编译、运行时和回归流程的整体收敛能力；如果团队既要维持开放生态接口，又要针对国产平台逐步形成本地优化和治理能力，那么混合式架构往往更值得投入，但前提是必须持续经营模块边界，而不是靠经验性补丁维持兼容。把这些判别准则讲清楚，比简单给出“先进性”结论更符合综述写作的目标。
